\documentclass[11pt, a4paper]{article}

% ── Packages ──────────────────────────────────────────────────────────────────
\usepackage{amsmath, amssymb}
\usepackage{booktabs}
\usepackage{enumitem}
\usepackage[T1]{fontenc}
\usepackage{geometry}
\usepackage{graphicx}
\usepackage[utf8]{inputenc}
\usepackage{listings}
\usepackage{microtype}
\usepackage{parskip}
\usepackage{titlesec}
\usepackage{titletoc}
\usepackage{xcolor}

% hyperref last
\usepackage[
  colorlinks = true,
  linkcolor  = teal!70!black,
  citecolor  = teal!70!black,
  urlcolor   = blue!60!black,
  pdftitle   = {Growing Metal Lace},
  pdfauthor  = {Carlo Perassi},
]{hyperref}

% ── Page geometry ─────────────────────────────────────────────────────────────
\geometry{margin=3cm, top=3.5cm, bottom=3.5cm}

% ── Colours ───────────────────────────────────────────────────────────────────
\definecolor{accent}{RGB}{0, 120, 120}
\definecolor{codebg}{RGB}{245, 247, 249}
\definecolor{codefg}{RGB}{40, 40, 60}

% ── Section formatting ────────────────────────────────────────────────────────
\titleformat{\section}
  {\large\bfseries\color{accent}}{{\thesection.}}{0.6em}{}[\color{accent}\titlerule]
\titleformat{\subsection}
  {\normalsize\bfseries\color{accent!80!black}}{{\thesubsection}}{0.5em}{}
\titleformat{\subsubsection}
  {\normalsize\itshape}{}{0em}{}

\titlespacing{\section}{0pt}{1.8em}{0.6em}
\titlespacing{\subsection}{0pt}{1.2em}{0.3em}

% ── TOC styling ───────────────────────────────────────────────────────────────
\titlecontents{section}[1.5em]
  {\smallskip\bfseries}
  {\contentslabel{1.5em}}{\hspace*{-1.5em}}
  {\titlerule*[6pt]{.}\contentspage}
\titlecontents{subsection}[3.5em]
  {}{\contentslabel{2em}}{}
  {\titlerule*[6pt]{.}\contentspage}

% ── Code listing ──────────────────────────────────────────────────────────────
\lstset{
  backgroundcolor = \color{codebg},
  basicstyle      = \small\ttfamily\color{codefg},
  breaklines      = true,
  frame           = single,
  framerule       = 0pt,
  rulecolor       = \color{codebg},
  xleftmargin     = 8pt,
  xrightmargin    = 8pt,
  aboveskip       = 8pt,
  belowskip       = 8pt,
  keywordstyle    = \bfseries\color{accent},
  commentstyle    = \itshape\color{gray},
  language        = Python,
}

% ── Title ─────────────────────────────────────────────────────────────────────
\title{%
  {\color{accent}\rule{\linewidth}{2pt}}\\[6pt]
  \textbf{\LARGE Growing Metal Lace}\\[4pt]
  \large From a Fractal Grammar to Evolutionary Geometry\\[2pt]
  \normalsize targeting DMLS/SLM Metal 3D Printing\\[6pt]
  {\color{accent}\rule{\linewidth}{1pt}}
}
\author{Carlo Perassi\\[4pt]
  \small\href{https://github.com/carlok/sierpsphere}%
             {\texttt{github.com/carlok/sierpsphere}}}
\date{2026}

% ── Commands ──────────────────────────────────────────────────────────────────
\newcommand{\code}[1]{\texttt{#1}}
\newcommand{\term}[1]{\emph{#1}}

% ─────────────────────────────────────────────────────────────────────────────
\begin{document}

\maketitle
\thispagestyle{empty}

% ── Abstract ──────────────────────────────────────────────────────────────────
\begin{abstract}
\noindent
We describe \textbf{SierpSphere}, a project that starts with a compact grammar
for fractal signed-distance functions (SDFs), layers a genetic algorithm (GA)
targeting real-world metal 3D-printing constraints
(DMLS/SLM~\cite{yap2015high}), and accelerates the inner evaluation loop on
Apple Silicon via PyTorch MPS~\cite{pytorch2019}.
The result is an automated search over a combinatorial space of printable,
aesthetically complex fractal objects — each one reconstructible from a single
line of symbolic notation — running at interactive speed on a laptop.
\end{abstract}

\vspace{1em}
{\small\tableofcontents}

\newpage

% ══ 1 ══════════════════════════════════════════════════════════════════════════
\section{Motivation}
\label{sec:motivation}

The Sierpi\'nski gasket~\cite{sierpinski1915} has been rendered millions of
times. What is less common is asking: \term{what other objects live near it in
the space of self-similar Boolean-SDF grammars?} And, once we can enumerate
that neighbourhood, \term{which of those objects could actually exist as a
centimetre-scale metal artefact?}

SierpSphere began as a weekend SDF renderer and grew into an end-to-end
pipeline spanning grammar design, real-time browser visualisation,
evolutionary search, and GPU acceleration — all within a single repository,
reproducible via Podman or a single shell script.

% ══ 2 ══════════════════════════════════════════════════════════════════════════
\section{The Grammar}
\label{sec:grammar}

\subsection{Signed Distance Functions as a generative substrate}
\label{subsec:sdf}

An SDF $f : \mathbb{R}^3 \to \mathbb{R}$ implicitly defines a surface as
$\{x \mid f(x) = 0\}$~\cite{quilez_sdf}.
Smooth Boolean operations — smooth-union, smooth-subtract, smooth-intersect
— compose SDFs without aliasing artifacts, and their blending parameter $k$
controls whether joins are sharp or organic (see \autoref{eq:smin}).

\begin{equation}
  \label{eq:smin}
  \mathrm{smin}(a,b,k) =
    \min(a,b) - k\,h(1-h), \quad
    h = \mathrm{clamp}\!\left(\tfrac{1}{2} + \tfrac{b-a}{2k},\,0,\,1\right)
\end{equation}

Polyhedral symmetry groups ($T_d$, $O_h$, $I_h$) provide a natural set of
axes along which to place child primitives~\cite{coxeter1973regular}.
At each iteration step the grammar places scaled copies of a primitive
(sphere, cube, or octahedron) along those axes and combines them with
the accumulated SDF via one of the three Boolean operations.

\subsection{A compact notation}
\label{subsec:notation}

A grammar is a symmetry group, a seed primitive, and an ordered list of steps.
The full state needed to reconstruct any evolved object fits on one line
(full formal definition in \href{notation.pdf}{\code{notation.tex}}):

\begin{equation}
  \label{eq:grammar}
  \Gamma = \mathsf{G}[\mathrm{P_0}] : s_1 \cdot s_2 \cdots s_n
\end{equation}

For example, the classic Sierpi\'nski sphere is
$\mathsf{T_d}[\mathrm{S}]:\ominus_s^{0.5}{}_{0.02}\cdot\oplus_s^{0.5}{}_{0.01}\cdot\ominus_s^{0.5}{}_{0.005}$.
A POSIX-safe filename slug encoding (replacing shell-unsafe characters with
letters) makes this notation directly usable in gallery filenames:
\code{Td.S.Ns4Ps2Ns1}.

\subsection{Mesh extraction}
\label{subsec:mesh}

The SDF is sampled on a uniform grid and converted to a triangular mesh via
the marching cubes algorithm~\cite{lorensen1987marching}.
At evaluation time a $28^3$ grid is used for speed; the epoch winner is
re-extracted at $64^3$ for the gallery GLB file.

% ══ 3 ══════════════════════════════════════════════════════════════════════════
\section{The Fitness Landscape}
\label{sec:fitness}

\subsection{Beauty as a measurable quantity}
\label{subsec:aesthetic}

We approximate visual complexity with metrics computed on the extracted mesh
(\autoref{tab:weights}):

\begin{itemize}[leftmargin=1.8em]
  \item \textbf{Fractal dimension} — box-counting on mesh vertices~\cite{mandelbrot1982fractal},
        rewarded when the slope $d_f \in [2.3, 2.7]$ (Gaussian centred at 2.5).
  \item \textbf{Curvature variance} — discrete mean curvature measure~\cite{meyer2003discrete};
        high variance implies rich surface detail.
  \item \textbf{Genus} — topological handles; a high-genus mesh has many
        through-holes, characteristic of Sierpi\'nski-like lace.
  \item \textbf{Normalised $S/V$ ratio} — rewards forms that fill space
        efficiently without being solid blobs.
  \item \textbf{Primitive diversity} — rewards grammars that mix sphere,
        cube, and octahedron across steps; sustains structural variety under
        selection pressure (see \autoref{subsec:diversity}).
\end{itemize}

\subsection{Manufacturability as a hard constraint}
\label{subsec:manufacturing}

Direct Metal Laser Sintering (DMLS) and Selective Laser Melting
(SLM)~\cite{yap2015high} impose strict geometric requirements.
Three hard gates eliminate any candidate that fails them outright:

\begin{enumerate}[leftmargin=1.8em]
  \item \textbf{No islands} — the mesh must be a single connected component.
  \item \textbf{No enclosed voids} — sealed internal cavities trap metal
        powder and cannot be removed post-print.
  \item \textbf{Minimum wall thickness $\geq 0.3\,\mathrm{mm}$} — estimated
        via ray-casting on the surface at the target scale of $80\,\mathrm{mm}$.
\end{enumerate}

Soft manufacturing metrics — overhang fraction, thermal mass uniformity,
drain openings — contribute to the weighted fitness score $f \in [0,1]$.

\subsection{Full fitness function}
\label{subsec:fitness_full}

Fifteen sub-scores are combined as a weighted sum (\autoref{tab:weights}).
Aesthetic metrics account for 53\% of the total weight;
manufacturing metrics for 47\%.
The two groups are deliberately balanced: a beautiful but unprintable object
scores zero, and a printable but trivial object scores poorly.

\begin{table}[h]
\centering
\caption{Fitness sub-scores and weights.}
\label{tab:weights}
\small
\begin{tabular}{llr}
\toprule
Metric & Group & Weight \\
\midrule
Fractal dimension       & Aesthetic       & 13\% \\
Min wall thickness      & Manufacturing   & 15\% \\
Curvature variance      & Aesthetic       &  9\% \\
Min feature size        & Manufacturing   &  8\% \\
Symmetry preservation   & Aesthetic       &  8\% \\
Drain openings          & Manufacturing   &  8\% \\
Normalised $S/V$        & Aesthetic       &  8\% \\
Support volume ratio    & Manufacturing   &  5\% \\
Thermal mass variance   & Manufacturing   &  5\% \\
Silhouette complexity   & Aesthetic       &  5\% \\
Genus                   & Aesthetic       &  5\% \\
Enclosed voids (gate)   & Manufacturing   &  3\% \\
No islands (gate)       & Manufacturing   &  3\% \\
Aspect ratio            & Aesthetic       &  3\% \\
Primitive diversity     & Aesthetic       &  2\% \\
\bottomrule
\end{tabular}
\end{table}

% ══ 4 ══════════════════════════════════════════════════════════════════════════
\section{The Genetic Algorithm}
\label{sec:ga}

\subsection{Overview}
\label{subsec:ga_overview}

A population of $N=24$ grammars~\cite{holland1975adaptation} evolves over 50
epochs. Each epoch:

\begin{enumerate}[leftmargin=1.8em]
  \item \textbf{Evaluate} — extract the mesh via marching cubes
        (\autoref{subsec:mesh}) and compute the fitness
        (\autoref{subsec:fitness_full}).
  \item \textbf{Select} — tournament selection ($k=2$, weak pressure) with
        global elitism (top 2 carried unchanged) plus \term{type-aware elitism}:
        the best viable individual of each seed type (sphere, cube, octahedron)
        is always preserved, preventing premature convergence to a single
        primitive family.
  \item \textbf{Reproduce} — single-point crossover on the iteration-step
        list (50\% probability) followed by per-field Gaussian mutation
        (rate 0.35 per step).
\end{enumerate}

\subsection{Mutation operators}
\label{subsec:mutation}

Numeric parameters (scale $\sigma$, distance $\delta$, smoothness $\rho$)
are jittered multiplicatively:
\[
  x' = \mathrm{clamp}\!\left(x + x \cdot r \cdot \mathcal{N}(0,1),\; x_{\min},\; x_{\max}\right)
\]
Categorical fields (operation type, primitive, seed type) are swapped by
uniform random draw.
Steps can be added or removed, allowing the grammar length to evolve within
$[1, 4]$.

\subsection{Diverse initial population}
\label{subsec:init}

The initial population is generated by \code{diverse\_population()}: equal
thirds of sphere-, cube-, and octahedron-seeded individuals, shuffled, with
step counts drawn uniformly from $[1,4]$ and all parameters fully randomised.
The first step always carves (subtract or intersect) to avoid the hard-gate
failure mode of pure-union blob grammars.

\subsection{Primitive diversity reward}
\label{subsec:diversity}

The \code{primitive\_diversity} sub-score (weight 2\%) rewards grammars that
mix sphere, cube, and octahedron across steps and seed.
A grammar using all three primitive types scores the full 0.02; two types
scores 0.01; one type scores 0.
A small bonus applies when step primitives vary independently of the seed.
This signal is lightweight but sufficient to break ties in favour of
structurally varied forms and to sustain diversity under selection pressure.

% ══ 5 ══════════════════════════════════════════════════════════════════════════
\section{GPU Acceleration via PyTorch MPS}
\label{sec:metal}

\subsection{The bottleneck}
\label{subsec:bottleneck}

An SDF grammar with $n$ operations evaluated on a $40^3$ grid requires
$40^3 \times n \approx 10^6$ floating-point evaluations per individual.
With 24 individuals and two CPU workers inside a Podman Linux VM on macOS,
epoch time reached 2 hours — too slow for iterative experimentation.

\subsection{Vectorised SDF on Metal}
\label{subsec:mps}

Apple Silicon exposes a unified-memory GPU via Metal Performance
Shaders~\cite{apple_mps}.
PyTorch's MPS backend~\cite{pytorch2019} maps tensor operations onto Metal
compute kernels transparently.

SDF evaluation is inherently data-parallel: all $N^3$ grid points are
independent. We reshape the grid into a flat $(N^3, 3)$ tensor and evaluate
all primitives and Boolean ops as batched arithmetic — no Python loop over
voxels:

\begin{lstlisting}[caption={Vectorised SDF evaluation on MPS (simplified).}]
pts = torch.stack([gx, gy, gz], dim=1)   # (N^3, 3) on MPS device
sdf = sd_sphere(pts, center, r)
for op in grammar_ops:
    child = sd_primitive(pts, op["type"], op["center"], op["radius"])
    sdf   = smooth_subtract(sdf, child, k=op["smooth_k"])
grid = sdf.reshape(N, N, N).cpu().numpy()
\end{lstlisting}

Marching cubes (scikit-image~\cite{virtanen2020scipy}) and trimesh fitness
metrics remain on CPU; only the SDF sampling moves to the GPU — sufficient
because SDF sampling dominates per-individual cost.

\subsection{Observed speedup}
\label{subsec:speedup}

On an Apple M-series chip with MPS, per-individual evaluation at $28^3$
drops from ${\sim}5\,\mathrm{min}$ (CPU, Podman VM) to ${\sim}3\,\mathrm{s}$
(Metal, native) — a speedup of roughly $100\times$.
Epoch time for $N=24$ falls from hours to minutes.

% ══ 6 ══════════════════════════════════════════════════════════════════════════
\section{System Architecture}
\label{sec:architecture}

The project is structured as five independently usable components
(\autoref{tab:components}):

\begin{table}[h]
\centering
\caption{Project components.}
\label{tab:components}
\begin{tabular}{lll}
\toprule
Component & Technology & Purpose \\
\midrule
\code{engine/}     & Python · Flask · trimesh  & SDF eval · GLB export · REST API \\
\code{viewer/}     & Three.js · GLSL · Nginx   & Real-time MC viewer · HQ PBR snapshot \\
\code{evolver/}    & Python · PyTorch MPS      & GA · fitness · Metal eval \\
\code{glb-viewer/} & Three.js (standalone)     & Material preview · turntable · 4K PNG \\
\code{tex/}        & \LaTeX                    & Grammar notation · this document \\
\bottomrule
\end{tabular}
\end{table}

The Podman/Docker stack (\code{docker-compose.yml}) orchestrates the
containerised path; \code{run\_native.sh} provides a one-command native macOS
setup with automatic virtualenv creation and dependency installation.

% ══ 7 ══════════════════════════════════════════════════════════════════════════
\section{Open Questions and Future Directions}
\label{sec:future}

\begin{itemize}[leftmargin=1.8em]
  \item \textbf{Higher-resolution output} — the epoch winner is saved at
        $64^3$; a post-evolution pass at $256^3$ on Metal would produce
        print-ready geometry in seconds.
  \item \textbf{Multi-objective Pareto front} — rather than a scalar fitness,
        maintain a Pareto front of (beauty, printability) pairs and let the
        user navigate the trade-off interactively~\cite{deb2002nsga}.
  \item \textbf{Meta-evolution} — the mutation operators themselves could be
        subject to evolution (hyper-heuristics~\cite{burke2013hyper}), allowing
        the search to adapt its own exploration strategy.
  \item \textbf{Differentiable SDF} — replacing marching cubes with a
        differentiable renderer (sphere tracing with PyTorch
        autograd~\cite{park2019deepsdf}) would enable gradient-based
        optimisation of grammar parameters.
  \item \textbf{Physical simulation} — coupling the fitness function to a
        finite-element stress simulator would reward geometries that are not
        just printable but structurally sound under load.
\end{itemize}

% ══ 8 (General audience) ══════════════════════════════════════════════════════
\section*{Short version (general audience / LinkedIn)}
\addcontentsline{toc}{section}{Short version (general audience)}

\textit{What happens when you let a computer invent its own fractal sculptures
and grade them on whether they could be 3D printed in metal?}

I built SierpSphere: a system where a grammar defines fractal shapes as
sequences of sphere-carving operations, a genetic algorithm evolves thousands
of variants, and a fitness function scores each one on visual complexity
\emph{and} real manufacturing constraints — minimum wall thickness, no
floating islands, no trapped powder.

The twist: evaluating one shape required sampling a million points in 3D space.
On a standard CPU that took minutes per shape. By rewriting the inner loop as
a PyTorch tensor operation on Apple Metal (MPS), the same evaluation runs in
about 3 seconds — a $100\times$ speedup without touching the algorithm or the
output quality.

The best shapes from each generation are saved as GLB files, and each one has
a compact symbolic name (like a chemical formula) that fully describes how to
recreate it.

\medskip
\noindent Source: \url{https://github.com/carlok/sierpsphere}

% ── References ────────────────────────────────────────────────────────────────
\newpage
\bibliographystyle{plain}

\begin{thebibliography}{99}

\bibitem{sierpinski1915}
W.~Sierpi\'nski.
\newblock Sur une courbe dont tout point est un point de ramification.
\newblock \emph{Comptes Rendus Acad. Sci. Paris}, 160:302--305, 1915.

\bibitem{quilez_sdf}
I.~Qui\~{n}onez (Inigo Quilez).
\newblock Signed distance functions.
\newblock \url{https://iquilezles.org/articles/distfunctions/}, 2008--2024.
\newblock Accessed 2026.

\bibitem{lorensen1987marching}
W.~E. Lorensen and H.~E. Cline.
\newblock Marching cubes: A high resolution 3D surface construction algorithm.
\newblock \emph{ACM SIGGRAPH Computer Graphics}, 21(4):163--169, 1987.

\bibitem{mandelbrot1982fractal}
B.~B. Mandelbrot.
\newblock \emph{The Fractal Geometry of Nature}.
\newblock W.H. Freeman, 1982.

\bibitem{coxeter1973regular}
H.~S.~M. Coxeter.
\newblock \emph{Regular Polytopes}.
\newblock Dover Publications, 3rd edition, 1973.

\bibitem{yap2015high}
C.~Y. Yap, C.~K. Chua, Z.~L. Dong, Z.~H. Liu, D.~Q. Zhang, L.~E. Loh, and
  S.~L. Sing.
\newblock Review of selective laser melting: Materials and applications.
\newblock \emph{Applied Physics Reviews}, 2(4):041101, 2015.

\bibitem{holland1975adaptation}
J.~H. Holland.
\newblock \emph{Adaptation in Natural and Artificial Systems}.
\newblock University of Michigan Press, 1975.

\bibitem{pytorch2019}
A.~Paszke et al.
\newblock {PyTorch}: An imperative style, high-performance deep learning
  library.
\newblock In \emph{Advances in Neural Information Processing Systems},
  volume~32, 2019.

\bibitem{apple_mps}
Apple Inc.
\newblock Metal performance shaders.
\newblock \url{https://developer.apple.com/documentation/metalperformanceshaders},
  2014--2024.

\bibitem{virtanen2020scipy}
P.~Virtanen et al.
\newblock {SciPy} 1.0: Fundamental algorithms for scientific computing in
  Python.
\newblock \emph{Nature Methods}, 17:261--272, 2020.

\bibitem{meyer2003discrete}
M.~Meyer, M.~Desbrun, P.~Schr\"{o}der, and A.~H. Barr.
\newblock Discrete differential-geometry operators for triangulated
  2-manifolds.
\newblock In \emph{Visualization and Mathematics III}, pages 35--57. Springer,
  2003.

\bibitem{deb2002nsga}
K.~Deb, A.~Pratap, S.~Agarwal, and T.~Meyarivan.
\newblock A fast and elitist multiobjective genetic algorithm: {NSGA-II}.
\newblock \emph{IEEE Transactions on Evolutionary Computation}, 6(2):182--197,
  2002.

\bibitem{burke2013hyper}
E.~K. Burke et al.
\newblock Hyper-heuristics: A survey of the state of the art.
\newblock \emph{Journal of the Operational Research Society},
  64(12):1695--1724, 2013.

\bibitem{park2019deepsdf}
J.~J. Park, P.~Florence, J.~Straub, R.~Newcombe, and S.~Lovegrove.
\newblock {DeepSDF}: Learning continuous signed distance functions for shape
  representation.
\newblock In \emph{CVPR}, pages 165--174, 2019.

\end{thebibliography}

\end{document}
